\subsubsection{Generic-fibre generation in arbitrary dimension}

\begin{proposition}[Proved: generation from the geometric generic fibre]
\label{D:generic-fibre-generation}
Let $k$ be algebraically closed, let $X$ be a normal connected
proper $k$-scheme, and let $f:X\to B$ be a dominant morphism
to a smooth connected proper curve.  Suppose the geometric
generic fibre $X_{\overline\eta}$ is integral.  Let
$Y\subsetneq X$ be closed and put $U=X\setminus Y$.
Then the following assertions hold.
\begin{enumerate}
\item If $\pi_1(U_{\overline\eta})$ is topologically finitely
generated and $U$ meets every fibre of $f$, then $\pi_1(U)$
is topologically finitely generated.
\item If $X_{\overline\eta}$ is normal and
$Q(X_{\overline\eta},Y_{\overline\eta})$ is topologically
finitely generated, then $Q(X,Y)$ is topologically finitely
generated.  In this assertion whole fibres may be contained
in $Y$.
\end{enumerate}
No properness assumption is imposed on $U\to B$, and the
generic fibre need not be smooth.
\end{proposition}
\begin{proof}
Put $F=k(B)$.  The arithmetic fundamental-group sequence
over $F$ is exact:
\[
 1\longrightarrow\pi_1(U_{\overline\eta})
  \longrightarrow\pi_1(U_\eta)
  \longrightarrow\operatorname{Gal}(F^{\mathrm{sep}}/F)
  \longrightarrow1;
\]
see \cite[Lemma 58.14.3, tag 0BTX]{C_StacksGenericFibre}.
Moreover, $\pi_1(U_\eta)\to\pi_1(U)$ is surjective.
Indeed, every connected finite \'etale cover of the normal
integral $U$ is integral, and its nonempty localization at
the generic point of the base remains integral.
Consequently the image of $\pi_1(U_{\overline\eta})$ in
\emph{any} quotient $G$ of $\pi_1(U)$ is closed and normal.
Denote this image by $N$, and write $Q=G/N$.  There is a
surjection $\operatorname{Gal}(F^{\mathrm{sep}}/F)\to Q$.

For assertion (1) we take $G=\pi_1(U)$; its subgroup $N$
is finitely generated by hypothesis.  For assertion (2)
we take $G=Q(X,Y)$.  In this case the map
$\pi_1(U_{\overline\eta})\to G$ factors through
$Q(X_{\overline\eta},Y_{\overline\eta})$, so $N$ is again
finitely generated.  To justify this factorization, every
codimension-one boundary point of $X_{\overline\eta}$ lies
above a horizontal codimension-one boundary point of $X$.
This follows from dimension, since the geometric base-field
extension is algebraic.  Both local rings are DVRs, the
first by the assumed geometric normality and the second
by normality of $X$.  Tameness of an original boundary
valuation persists under this extension of DVRs, by
\cite[Lemma 15.116.9, tag 0EXY]{C_StacksLogFibration}.
Thus every finite cover counted by $G$ restricts to a
cover counted by the indicated geometric fixed-boundary
group.

We next show that $Q$ comes from one fixed open curve.
The morphism $U\to B$ has a nonempty smooth locus dominating
$B$.  Indeed its generic fibre is geometrically integral,
hence contains a dense smooth open by
\cite[Lemma 33.25.7, tag 056V]{C_StacksGenericFibre}.
The smooth locus is open and a smooth morphism is open.
Its image therefore contains a nonempty open
$B^\circ\subset B$.  In particular, for every closed
$b\in B^\circ$, there is a retained fibre component of
multiplicity one: take a component through an $f$-smooth
point of $U$ above $b$.

Let $Q_0$ be a finite quotient of $Q$.  The corresponding
connected cover $V\to U$ becomes over the generic fibre
the pullback of a connected finite Galois extension
$F'/F$ with group $Q_0$, by the arithmetic exact sequence.
Thus $k(V)=k(X)F'$.  Normalize $B$ in $F'$ to get
$B'\to B$.  If $E$ is a component of the fibre over $b$
with multiplicity $m_E$, and $e_E$ is the ramification
index of any prolongation of its valuation to $k(V)$,
multiplication of value-group indices gives
\[
                         e(b'/b)\mid m_Ee_E.
\tag{D.GF.1}
\]
At a retained component $e_E=1$, because $V\to U$ is
\'etale.  Therefore $e(b'/b)=1$ for $b\in B^\circ$.
The residue field at every closed point of $B'$ is $k$,
so $B'\to B$ is unramified, hence \'etale, over $B^\circ$.
This holds for every finite quotient of $Q$.  Consequently
the surjection from the absolute Galois group factors as
\[
                         \pi_1(B^\circ)\longrightarrow Q.
\]

In assertion (1), $U$ meets every fibre.  Hence above each
of the finitely many omitted $b$, some fibre component
has its generic point in $U$.  The index (D.GF.1) implies
that the order of the image of the full inertia at $b$
in every finite quotient of $Q$ divides that component's
multiplicity.  Thus the image in $Q$ of the full inertia
is finite.  Choose an open normal subgroup $Q'\subset Q$
meeting these finitely many finite inertia images trivially.
The inverse image of $Q'$ corresponds to a finite \'etale
cover $B_1^\circ\to B^\circ$.  After its smooth proper
compactification $B_1$, the map
$\pi_1(B_1^\circ)\to Q'$ kills all boundary inertia and
factors through $\pi_1(B_1)$.  This group is finitely
generated, so $Q'$, $Q$, and then $G$ are finitely generated.

In assertion (2), a fibre component is either retained,
in which case $e_E=1$, or is an original boundary divisor,
in which case $e_E$ is prime to the positive characteristic
$p$.  Equation (D.GF.1) therefore bounds the $p$-part of
$e(b'/b)$ by the $p$-part of $m_E$ at every component.
The wild inertia image at each omitted $b$ in $Q$ is
finite.  In characteristic zero it is trivial.  Choose
an open normal subgroup $Q'$ avoiding these finite wild
images.  For the corresponding open curve $B_1^\circ$,
every wild inertia group maps into a conjugate of an
original wild inertia group and hence trivially to $Q'$.
Thus $Q'$ is a quotient of $\pi_1^{\mathrm t}(B_1^\circ)$,
which is finitely generated.  This proves finite generation
of $Q$ and $G$ in assertion (2).

The finiteness assertion about a profinite inertia image
uses only its bounded finite quotient orders, as in
Proposition~\ref{D:vertical-fibration-multiplicity}; the
existence of one open normal subgroup avoiding finitely
many finite images follows by separating their nonidentity
elements in finite quotients.  The proper and tame curve
generation theorems used in the last two paragraphs are
\cite[X, Theorem 2.6; XIII, Corollary 2.12]{SGA1}.
\end{proof}

\begin{corollary}[Proved: arbitrary-dimensional vertical boundary]
\label{D:vertical-boundary-all-dimensions}
Let $X$ be a normal connected proper scheme over an
algebraically closed field and let $f:X\to B$ be a
nonconstant morphism to a proper curve.  Let
$Y\subsetneq X$ be a closed subset with $f(Y)$ finite.
Then $Q(X,Y)$ is topologically finitely generated.
If, after Stein factorization of $f$, the open $X\setminus Y$
meets every fibre, its full fundamental group is
topologically finitely generated.  The latter condition
is automatic if $Y$ has codimension at least two in $X$.
\end{corollary}
\begin{proof}
Take Stein factorization, so the new base is a smooth
connected proper curve, $f_*\mathcal O_X=\mathcal O_B$,
and $F=k(B)$ is algebraically closed in $k(X)$.
The p-rank-one argument in
Theorem~\ref{D:surface-map-to-curve} is independent of
$\dim X$: $[F:F^p]=p$, and a non-$p$th-power of $F$
cannot acquire a $p$th root in $k(X)$ without contradicting
relative algebraic closedness.  Thus $k(X)/F$ is regular
and the geometric generic fibre is integral.

Since the boundary is vertical, $U_{\overline\eta}=
X_{\overline\eta}$ is proper and connected.  Its fundamental
group is finitely generated by
\cite[X, Theorem 2.9, pp. 273--274]{SGA1}, with no
normality or smoothness requirement.  For the full group
under the stated fibre-intersection hypothesis, assertion
(1) of Proposition~\ref{D:generic-fibre-generation} applies.
For $Q(X,Y)$ the same proof applies verbatim: its geometric
generic image is a quotient of this proper fundamental
group, and all its finite covers are tame at every deleted
fibre component.  Consequently (D.GF.1) bounds the wild
base inertia and proves finite generation without imposing
geometric normality on the generic fibre.

Finally, every fibre of a dominant morphism from a normal
integral scheme to a smooth curve has a codimension-one
component: a local uniformizer cuts a nonzero principal
divisor.  A subset of codimension at least two cannot
contain such a component.  Thus its complement meets
every fibre.
\end{proof}

\begin{corollary}[Proved: projective normalization of the geometric generic fibre]
\label{D:projective-generic-normalization}
Let $X$ be a normal connected proper scheme over an
algebraically closed field.  Suppose $X$ admits a
nonconstant morphism to a proper curve such that, after
Stein factorization, the normalization of its geometric
generic fibre is projective.  Then $Q(X,Y)$ is
topologically finitely generated for every proper closed
subset $Y\subset X$.  In particular, when $Y$ has
codimension at least two, the full group
$\pi_1(X\setminus Y)$ is topologically finitely generated.
\end{corollary}
\begin{proof}
The same Stein and p-rank-one arguments make the generic
fibre geometrically integral.  By
\cite[Lemma 33.27.4, tag 0BXT]{C_StacksLogFibration}, after
a finite purely inseparable extension of the base function
field its normalization is geometrically normal.  Normalize
the base curve in this extension and normalize $X$ in the
compositum of function fields, obtaining $X'\to B'$.
As in Theorem~\ref{D:surface-map-all-boundaries}, the map
$X'\to X$ is finite and radicial: a fixed $p$-power of
every integral element belongs to the original normal
ring.  In characteristic zero this step is unnecessary.

The geometric generic fibre of $X'$ is the normalization
of the original geometric generic fibre.  Indeed, its
finite birational map to that fibre has normal source,
so it has the defining property of normalization.  It is
therefore projective by hypothesis.  Proposition
\ref{D:fixed-boundary-positive} gives finite generation
of its fixed-boundary tame group for the inverse image
of $Y$.  Assertion (2) of
Proposition~\ref{D:generic-fibre-generation} proves that
$Q(X',Y')$ is finitely generated.

The finite radicial map preserves the \'etale fundamental
group of the open, and each new codimension-one boundary
DVR lies over an original one.  Tame DVR base change
therefore gives a surjection $Q(X',Y')\to Q(X,Y)$,
exactly as in Theorem~\ref{D:surface-map-all-boundaries}.
This proves the assertion.  If the original boundary
has codimension at least two, no wild inertia is killed
in defining $Q(X,Y)$, so it equals the full open
fundamental group.
\end{proof}
